\documentclass[utf8]{article} %中文文档类
\usepackage[left=2.50cm, right=2.50cm, top=2.50cm, bottom=2.50cm]{geometry} % 页边距
\usepackage{indentfirst} % 首行缩进
\usepackage{fancyhdr,fancybox} % 设置页眉、页脚
\usepackage{ctexcap} % 标题是中文的
\usepackage{helvet} % 用来指定beamer使用的字体
\usepackage{hyperref} % bookmarks 
\usepackage{multicol} % 分栏
\usepackage{cite} % 引用

\usepackage{graphicx} % 图片
\usepackage{subcaption} % 子图
\usepackage{multirow} % 表格
\usepackage{float} % 浮动体
\usepackage{booktabs} % 三线表
\usepackage{longtable} % 长表格
\usepackage{diagbox} % 表头斜线
\usepackage{multirow} % 表格合并
\usepackage{makecell} % 表格换行
\usepackage{enumitem} % 枚举环境宏包
\usepackage{framed} % 边框包

\usepackage{amsmath, amsfonts, amssymb} % 数学公式、符号
\numberwithin{equation}{section}   % 编号从一级标题开始
\usepackage[english]{babel} % 数学公式标准
\usepackage{bm} % 加粗方程字体 
\usepackage{caption} % 标题
\usepackage{color} % 配色
\usepackage[braket]{qcircuit} % 启用狄拉克符号支持, 用于绘制电路图
\usepackage{physics} % 量子力学中的符号加载

\newtheorem{Def}{\hskip 2em {\color{green} 定义}}[subsection]
\newtheorem{Thm}{\hskip 2em {\color{blue} 定理}}[subsection]
\newtheorem{Lem}{\hskip 2em {\color{magenta} 引理}}[subsection]
\newtheorem{Prop}{\hskip 2em {\color{cyan} 命题}}[subsection]
\newtheorem{Note}{\hskip 2em {\color{red} 注}}[subsection]
\newtheorem{Cor}{\hskip 2em {\color{cyan} 推论}}[subsection]
\newtheorem{Eg}{\hskip 2em {\color{black} 例}}[subsection]
\newenvironment{Proof}{{\noindent\it Proof:}\quad}{\hfill $\sharp $\par}
\allowdisplaybreaks

% 页眉页脚设置
\usepackage{fancyhdr}
\usepackage{lastpage} % 获得总页数
\pagestyle{fancy}
\fancyfoot{}
\lhead{\textit{Core Quantum Algorithm Components}}\chead{}\rhead{\textit{\thepage\ of \pageref{LastPage}}}
\lfoot{}\cfoot{}\rfoot{}
\renewcommand{\headrulewidth}{0.4pt}

% 标题设置
\title{\textbf{Module2 核心量子算法构件}}
\date{\today}
\author{\textbf{STU\quad 牛晓铭}}

\hypersetup{
	colorlinks=true,
	linkcolor=black
} % 使目录为黑色字
\begin{document}  
\maketitle
\renewcommand{\contentsname}{} 
\tableofcontents
\thispagestyle{empty}
	
\newpage
\setcounter{page}{1}
\section{引言: 量子算法原语}

在上一模块中我们建立了量子计算的语言: 量子比特、酉门、电路与测量. 本模块介绍四类最基础的
\textbf{量子算法原语(quantum primitives)}: 量子傅里叶变换(QFT)、量子相位估计(QPE)、
振幅放大(AA)与哈密顿量模拟(Hamiltonian simulation, 含 Lie--Trotter 公式).
它们之所以被称为"构件", 是因为几乎所有量子算法都以它们为子程序:
QFT 是 QPE 的核心组成部分; QPE 是 Shor 分解、HHL 线性方程组、本征值问题、振幅估计的共同基础;
振幅放大在一切需要"放大成功概率"的场景中被反复调用; 哈密顿量模拟则通常作为 QPE 中酉算子
$U = e^{-iHt}$ 的实现方式被调用. 这些构件之间相互嵌套的关系构成了现代量子算法的骨架.

\begin{Note}
    按照 Lin Lin 讲义的划分, 上述四者(连同 Grover 搜索)构成最没有争议的量子原语集合.
    本模块先介绍不依赖其它构件的 QFT 与 AA, 再介绍把它们组合起来的 QPE 与 Trotter 型
    哈密顿量模拟. 更现代的构件(block-encoding、qubitization、QSVT)将在下一模块中
    统一并推广这些方法.
\end{Note}

\section{量子傅里叶变换 (QFT)}

\subsection{定义与基本性质}

\begin{Def}[QFT 的定义]
    设 $N = 2^n$. \textbf{量子傅里叶变换(quantum Fourier transform, QFT)}是
    $\mathbb{C}^N$ 上的酉算子 $U_{\mathrm{FT}}$, 作用在计算基态上为
    \begin{equation}
        U_{\mathrm{FT}} \ket{j} = \frac{1}{\sqrt{N}} \sum_{k \in [N]} e^{2\pi i\, jk/N}\, \ket{k},
        \qquad j \in [N].
    \end{equation}
    其逆变换为
    \begin{equation}
        U_{\mathrm{FT}}^\dagger \ket{j} = \frac{1}{\sqrt{N}} \sum_{k \in [N]} e^{-2\pi i\, jk/N}\, \ket{k}.
    \end{equation}
\end{Def}

\begin{Prop}[基本性质]
    \begin{enumerate}
        \item $U_{\mathrm{FT}}$ 是酉算子($U_{\mathrm{FT}}^\dagger U_{\mathrm{FT}} = I$);
        \item $U_{\mathrm{FT}} \ket{0^n} = \frac{1}{\sqrt{N}} \sum_{k} \ket{k} = H^{\otimes n}\ket{0^n}$,
        即 QFT 把全零态变为均匀叠加态;
        \item 与经典 FFT 的关键区别: QFT 的加速是指数级的(经典 FFT 需要 $O(N\log N)$ 次运算,
        而 QFT 只需要 $O(n^2) = O(\log^2 N)$ 个门), 但 QFT 的\textbf{输出是叠加态},
        测量后只能得到单个基态, 无法直接读出全部傅里叶系数 —— 必须通过后续算法
        (如相位估计)来利用这一变换.
    \end{enumerate}
\end{Prop}

\subsection{二进制分解与张量积结构}

把 $j, k$ 写成二进制 $j = (j_{n-1}\cdots j_0)$, $k = (k_{n-1}\cdots k_0)$, 则

\begin{equation}
    \frac{jk}{N} = k_0\, (0.j_{n-1}\cdots j_0) + k_1\, (0.j_{n-2}\cdots j_0) + \cdots + k_{n-1}\, (0.j_0),
\end{equation}

其中 $0.j_{m}\cdots j_0$ 表示二进制小数 $j_m/2 + \cdots + j_0/2^{m+1}$. 于是指数因子可以分解,
代入定义并逐比特求和, 得到 QFT 的张量积形式:

\begin{equation}
    U_{\mathrm{FT}} \ket{j_{n-1}\cdots j_0}
    = \frac{1}{\sqrt{2^n}} \Big(\ket{0} + e^{2\pi i (0.j_0)}\ket{1}\Big)
    \Big(\ket{0} + e^{2\pi i (0.j_1 j_0)}\ket{1}\Big)
    \cdots \Big(\ket{0} + e^{2\pi i (0.j_{n-1}\cdots j_0)}\ket{1}\Big).
\end{equation}

注意: 第一个因子(对应 $\ket{k_{n-1}}$)只依赖最低比特 $j_0$, 而最后一个因子
(对应 $\ket{k_0}$)依赖全部比特 —— 即\textbf{输出比特的位序与输入相反}.

\subsection{电路实现}

实现的关键是"受控相位旋转". 记单比特旋转门
\begin{equation}
    R_z(\varphi) = \begin{pmatrix} 1 & 0 \\ 0 & e^{i\varphi} \end{pmatrix},
    \qquad R_m := R_z(2\pi/2^{m}).
\end{equation}

\begin{Eg}[受控旋转的实现]
    目标: 把 $j_{n-1}, \dots, j_0$ 的相位信息收集到辅助比特上:
    \[ \ket{0}\ket{j} \;\longrightarrow\; \frac{1}{\sqrt{2}}\Big(\ket{0} + e^{2\pi i(0.j_{n-1}\cdots j_0)}\ket{1}\Big)\ket{j}. \]
    利用
    \[ e^{2\pi i (0.j_{n-1}\cdots j_0)} = e^{2\pi i (0.j_{n-1})}\, e^{2\pi i (0.0 j_{n-2})} \cdots e^{2\pi i (0.\underbrace{0\cdots 0}_{n-2}\, j_0)}, \]
    该操作等价于依次施加 $H$ 与以 $j_{m}$ 为控制比特的受控 $R_{m+1}$ 门
    ($m = 0, \dots, n-1$): 每个比特 $j_m$ 贡献相位 $e^{2\pi i (0.0\cdots 0 j_m)}$.
\end{Eg}

QFT 电路(以 3 比特为例, 输出位序已反转, 需用 SWAP 门恢复):

\begin{equation*}
\Qcircuit @C=.9em @R=.8em {
    \lstick{\ket{j_2}} & \gate{H} & \ctrl{1} & \ctrl{2} & \qw & \qw & \qw & \qw & \qw & \qw & \qw & \rstick{\ket{k_0}} \qw \\
    \lstick{\ket{j_1}} & \qw & \gate{R_2} & \qw & \gate{H} & \ctrl{1} & \qw & \qw & \qw & \qw & \qw & \rstick{\ket{k_1}} \qw \\
    \lstick{\ket{j_0}} & \qw & \qw & \gate{R_3} & \qw & \gate{R_2} & \gate{H} & \qw & \qw & \qw & \qw & \rstick{\ket{k_2}} \qw
}
\end{equation*}

\begin{Thm}[QFT 的复杂度]
    $n$ 比特 QFT 可以用 $O(n^2)$ 个单比特与受控门实现, 电路深度为 $O(n)$.
    逆变换 $U_{\mathrm{FT}}^\dagger$ 只需把所有受控旋转替换为共轭 $R_m^\dagger$ 并反向执行.
\end{Thm}

\section{量子相位估计 (QPE)}

\subsection{问题设定与 Hadamard 检验}

\begin{Def}[相位估计问题]
    设 $U$ 是酉算子, $\ket{\psi}$ 是其特征向量:
    \begin{equation}
        U \ket{\psi} = e^{2\pi i \theta} \ket{\psi}, \qquad \theta \in [0, 1).
    \end{equation}
    给定 oracle 访问 $U$(及其幂), 目标是估计 $\theta$ 到一定精度.
\end{Def}

\begin{Note}
    相位估计是最重要的量子原语之一: Shor 分解、HHL 线性方程组、本征值问题、振幅估计、
    量子行走等都以它为子程序. 经典做法 $j|U\ket{\psi} / j|\psi = e^{2\pi i\theta}$
    (逐元素除法)在量子计算机上无法高效实现, 需要新的方法.
\end{Note}

最简单的相位估计是 \textbf{Hadamard 检验}(已在 Module1 中介绍): 对控制比特先施加
$H$, 作用受控 $U$ 后再施加 $H$, 则测量结果为 $1$ 的概率满足

\begin{equation}
    p(1) = \frac{1}{2}\big(1 - \cos(2\pi\theta)\big),
\end{equation}

从而 $\theta = \frac{1}{2\pi}\arccos\big(1 - 2p(1)\big) \pmod 1$.
但该方法的精度受限于采样误差: 由 Module1 的分析, 达到加性精度 $\varepsilon$
需要 $O(1/\varepsilon^2)$ 次重复测量. 相位估计的任务正是把精度提高到 $O(1/\varepsilon)$.

\subsection{受控 $U^{2^j}$ 与二进制分解}

把相位 $\theta = (0.\theta_{d-1}\cdots\theta_0)$ 写成 $d$ 位二进制小数, 考虑一个"大"受控酉算子

\begin{equation}
    \mathcal{U} := \sum_{j \in [2^d]} \ket{j}\bra{j} \otimes U^j.
\end{equation}

看似需要实现全部 $2^d$ 个不同的 $U^j$, 但利用二进制展开
$j = \sum_i j_i 2^i$ 得 $U^j = \prod_i U^{j_i 2^i}$, 于是

\begin{equation}
    \mathcal{U} = \prod_{i=0}^{d-1} \Big(\ket{0}\bra{0} \otimes I + \ket{1}\bra{1} \otimes U^{2^i}\Big),
\end{equation}

即只需 $d$ 个受控 $U^{2^i}$ 门. 若 $U$ 可以被 \textbf{快速前推(fast-forwarded)}
(实现 $U^j$ 的代价与 $j$ 无关, 例如 $U = R_z(\varphi)$), 该分解比直接实现 $\mathcal{U}$
指数级节省资源.

\subsection{QFT 基 QPE 电路}

设 $T = 2^t$, 其中 $t$ 为相位寄存器比特数. QPE 的电路结构为:
\begin{equation*}
\Qcircuit @C=1em @R=.8em {
    \lstick{\ket{0}^{\otimes t}} & \gate{H^{\otimes t}} & \multigate{1}{\mathcal{U}} & \gate{U_{\mathrm{FT}}^\dagger} & \meter & \rstick{\tilde{\theta}} \qw \\
    \lstick{\ket{\psi}} & \qw & \ghost{\mathcal{U}} & \qw & \qw & \qw
}
\end{equation*}

推导如下:

\begin{equation}
    \ket{0^t}\ket{\psi} \;\xrightarrow{H^{\otimes t}}\;
    \frac{1}{\sqrt{T}}\sum_{j\in[T]} \ket{j}\ket{\psi}
    \;\xrightarrow{\mathcal{U}}\;
    \frac{1}{\sqrt{T}}\sum_{j\in[T]} e^{2\pi i j\theta} \ket{j}\ket{\psi}
    \;\xrightarrow{U_{\mathrm{FT}}^\dagger}\;
    \sum_{k\in[T]} \alpha_k \ket{k}\ket{\psi},
\end{equation}

其中

\begin{equation}
    \alpha_k = \frac{1}{T}\sum_{j\in[T]} e^{2\pi i j(\theta - k/T)}
    = \frac{1}{T}\cdot\frac{1 - e^{2\pi i T(\theta - k/T)}}{1 - e^{2\pi i (\theta - k/T)}}.
\end{equation}

\begin{Thm}[精确相位的确定性]
    若 $\theta = k_0/T$ 恰有 $t$ 位二进制表示(即 $\theta = (0.\theta_{t-1}\cdots\theta_0)$),
    则 $\alpha_k = \delta_{k,k_0}$, 测量相位寄存器\textbf{确定性地}得到 $k_0$.
\end{Thm}

\begin{Eg}[d = 1 的特殊情形]
    当 $t = 1$ 时 $U_{\mathrm{FT}} = H$, 上述电路退化为 Hadamard 检验.
    这说明 Hadamard 检验正是 QPE 在 1 比特相位寄存器下的特例;
    且此时 $\theta$ 不必是单比特数 —— 一般情形由下述误差分析处理.
\end{Eg}

\subsection{误差分析}

当 $\theta_0$ 不能由 $t$ 位二进制精确表示时, 测量结果是随机的.
记周期距离 $\norm{x}_1 := \min\{x \bmod 1,\ 1 - (x \bmod 1)\}$, 由几何级数公式
$|1 - e^{i\varphi}| \ge 2|\varphi|/\pi$ 可得

\begin{equation}
    |\alpha_k| \le \frac{1}{2T\, \norm{\theta_0 - k/T}_1}.
\end{equation}

\begin{Thm}[QPE 的精度与成功概率]
    设精度参数 $\varepsilon = 2^{-d}$, 相位寄存器比特数取
    \begin{equation}
        t = d + \lceil \log_2(\varepsilon^{-1}) \rceil,
    \end{equation}
    则最大模拟时间 $T = 2^t = O(\varepsilon^{-1})$, 且
    \begin{equation}
        P\Big(\norm{\theta_0 - k_0/T}_1 \le \varepsilon\Big) \ge 1 - \varepsilon,
    \end{equation}
    即测量得到的 $k_0$ 以概率 $\ge 1 - \varepsilon$ 给出 $\theta_0$ 的 $\varepsilon$ 近似.
    总查询复杂度为 $O(T) = O(1/\varepsilon)$, 比 Hadamard 检验的 $O(1/\varepsilon^2)$ 平方级改进.
\end{Thm}

\begin{Proof}
    由 $|\alpha_k| \le 1/(2T \norm{\theta_0 - k/T}_1)$, 对距离 $\theta_0$ 超过 $\varepsilon$ 的
    $k$ 求和并用积分放缩:
    \[ P \le \sum_{\norm{\theta_0 - k/T}_1 \ge \varepsilon} \frac{1}{4T^2 \norm{\theta_0 - k/T}_1^2}
       \le \frac{1}{2T\varepsilon} + \frac{1}{2(T\varepsilon)^2}. \]
    取 $T\varepsilon \ge 1$ 即得失败概率 $\le \varepsilon$ 量级. 这也是 Nielsen--Chuang
    书中结论 $t = n + \lceil\log_2(2 + 1/(2\varepsilon))\rceil$ 的同一尺度.
\end{Proof}

\begin{Note}[量子中位数法与两个假设的放松]
    上述分析默认 $\ket{\psi}$ 恰为特征向量且相位有精确二进制表示.
    (1) 若初态是特征向量的线性组合 $\ket{\psi} = \sum_k c_k \ket{k}$,
    则 QPE 以概率 $|c_0|^2$ 同时返回目标相位与本征态 —— 这正是基态能量估计和 HHL 的工作方式;
    (2) 为避免因 $\log_2(1/\varepsilon)$ 个额外比特增加电路深度, 可把 QPE 运行 $\log(1/\varepsilon)$ 次
    并取相位的中位数(\textbf{量子中位数法}), 以 $\log(1/\varepsilon)$ 次重复代替额外的比特数与模拟时间.
\end{Note}

\subsection{应用: 振幅估计}

振幅放大算子(下一节)的特征值携带成功概率的信息, 因此可以用 QPE 估计概率 ——
这是 QPE 与振幅放大结合的第一个重要应用.

\begin{Prop}[振幅估计]
    设 Grover 算子 $G$ 在二维不变子空间 $\mathrm{span}\{\ket{\mathrm{good}}, \ket{\mathrm{bad}}\}$ 上
    的矩阵表示为旋转角 $\theta$ 的旋转(即 $e^{-i\theta Y}$),
    其两个特征态为 $\ket{\pm} = \frac{1}{\sqrt{2}}(\ket{\mathrm{good}} \pm i\ket{\mathrm{bad}})$,
    对应特征值 $e^{\pm i\theta}$.
    用 QPE 估计 $\theta$ 到加性精度 $\varepsilon$, 即可得到 $p_0 = \sin^2(\theta/2)$ 的估计,
    总复杂度 $O(1/\varepsilon)$; 而直接蒙特卡洛采样需要 $O(1/\varepsilon^2)$ 次.
    若 $p_0$ 很小需乘性精度, 复杂度为 $O(p_0^{-1/2}\varepsilon^{-1})$, 相对采样的
    $O(p_0^{-1}\varepsilon^{-2})$ 仍是平方级改进.
\end{Prop}

\section{振幅放大 (AA) 与最优振幅放大 (OAA)}

\subsection{Grover 搜索: 振幅放大的原型}

\begin{Def}[无结构搜索问题]
    给定布尔函数 $f: \{0,1\}^n \to \{0,1\}$ 与 promise: 存在唯一标记态 $x_0$ 使 $f(x_0) = 1$,
    找出 $x_0$. 经典算法最坏情形需要 $N - 1 = 2^n - 1$ 次查询,
    量子算法可以做到 $O(\sqrt{N})$ 次 —— 这就是 Grover 的\textbf{二次加速}.
\end{Def}

设 oracle $U_f\ket{x}\ket{y} = \ket{x}\ket{y \oplus f(x)}$. 取 $\ket{y} = \ket{-}$,
由相位反冲得 $U_f\ket{x}\ket{-} = (-1)^{f(x)}\ket{x}\ket{-}$, 于是相位 oracle
$U_f\ket{x} = (-1)^{f(x)}\ket{x}$ 等价于\textbf{反射算子}

\begin{equation}
    R_{x_0} = I - 2\ket{x_0}\bra{x_0}.
\end{equation}

初态取均匀叠加 $\ket{\phi_0} = H^{\otimes n}\ket{0^n}$, 定义关于 $\ket{\phi_0}$ 的反射

\begin{equation}
    R_0 = 2\ket{\phi_0}\bra{\phi_0} - I = H^{\otimes n}\big(2\ket{0^n}\bra{0^n} - I\big)H^{\otimes n}.
\end{equation}

$R_0$ 可以由 $n$ 比特受控 $X$ 门与辅助比特 $\ket{-}$ 实现.

\begin{Thm}[Grover 迭代]
    记 $\theta = 2\arcsin(1/\sqrt{N})$, 分解
    $\ket{\phi_0} = \sin(\theta/2)\ket{x_0} + \cos(\theta/2)\ket{x_0^\perp}$.
    \textbf{Grover 算子} $G = R_0 R_{x_0}$ 在二维不变子空间
    $\mathrm{span}\{\ket{x_0}, \ket{x_0^\perp}\}$ 内是旋转角为 $\theta$ 的旋转:
    \begin{equation}
        [G] = \begin{pmatrix} \cos\theta & \sin\theta \\ -\sin\theta & \cos\theta \end{pmatrix},
    \end{equation}
    且
    \begin{equation}
        G^k \ket{\phi_0} = \sin\!\big((2k+1)\theta/2\big)\ket{x_0}
        + \cos\!\big((2k+1)\theta/2\big)\ket{x_0^\perp}.
    \end{equation}
\end{Thm}

\begin{Proof}
    两个反射的乘积是旋转, 且 $\mathrm{span}\{\ket{x_0}, \ket{x_0^\perp}\}$ 是不变子空间:
    $R_{x_0}$ 保持 $\ket{x_0^\perp}$ 并翻转 $\ket{x_0}$ 的符号, $R_0$ 保持 $\ket{\phi_0}$ 并翻转
    其正交方向. 直接计算 $G\ket{\phi_0} = \sin(3\theta/2)\ket{x_0} + \cos(3\theta/2)\ket{x_0^\perp}$
    可知角度从 $\theta/2$ 增为 $3\theta/2$, 每次迭代增加 $\theta$.
\end{Proof}

\begin{Cor}[查询复杂度]
    取 $k = O\big(1/\theta\big) = O(\sqrt{N})$ 次迭代, 即可使 $\ket{x_0}$ 的振幅达到
    $\Omega(1)$. 具体地, 测量失败概率不超过 $4/N$, 且输出可以经典验证
    ($f(x)$ 检查), 重复即可. 查询复杂度 $O(\sqrt{N})$ 与经典 $O(N)$ 相比是二次加速;
    且这是最优的 —— 任何量子算法至少需要 $\Omega(\sqrt{N})$ 次查询(见讲义 2.4 节的混合论证证明).
    若存在 $M$ 个标记态, 复杂度为 $O(\sqrt{N/M})$.
\end{Cor}

\subsection{振幅放大: 一般框架}

Grover 搜索只是振幅放大在"标记态 = 单个计算基态"时的特例.

\begin{Thm}[振幅放大] \label{thm:AA}
    设 $\Pi_{\mathrm{good}}$ 是"好空间"上的投影, 初态由 oracle $U_0$ 制备:
    \begin{equation}
        \ket{\phi_0} = U_0 \ket{0^n}, \qquad
        \ket{\phi_0} = \Pi_{\mathrm{good}}\ket{\phi_0} + (I - \Pi_{\mathrm{good}})\ket{\phi_0}
        = \sqrt{p}\,\ket{\mathrm{good}} + \sqrt{1-p}\,\ket{\mathrm{bad}},
    \end{equation}
    其中 $p = \bra{\phi_0}\Pi_{\mathrm{good}}\ket{\phi_0}$, $\ket{\mathrm{good}}$,
    $\ket{\mathrm{bad}}$ 正交, 记 $p = \sin^2(\theta/2)$.
    假设可以访问反射 $R_{\mathrm{good}} = I - 2\Pi_{\mathrm{good}}$ 与
    $R_0 = 2\ket{\phi_0}\bra{\phi_0} - I = U_0(2\ket{0^n}\bra{0^n} - I)U_0^\dagger$,
    定义 $G = R_0 R_{\mathrm{good}}$, 则对任意 $k \ge 0$,
    \begin{equation}
        \big|\bra{\mathrm{good}} G^k \ket{\phi_0}\big| = \sin\!\big((2k+1)\theta/2\big).
    \end{equation}
    于是 $k = O(1/\theta) = O(1/\sqrt{p})$ 次迭代后, 成功概率达到 $\Omega(1)$.
\end{Thm}

\begin{Eg}[信号比特场景]
    若 $U_0$ 需要 $m$ 个辅助(信号)比特, 即
    \[ U_0 \ket{0^m}\ket{0^n} = \sqrt{p}\,\ket{0^m}\ket{\phi} + \sqrt{1-p}\,\ket{\phi^\perp}, \]
    其中 $\ket{\phi^\perp}$ 与 $\ket{0^m}$ 正交. 此时"好"的状态由辅助比特全 $0$ 标志,
    反射简化为 $R_{\mathrm{good}} = (I - 2\ket{0^m}\bra{0^m}) \otimes I$, 后选择测量辅助比特
    即得到成功概率 $p$. 这类"可验证成功"的结构在 block-encoding 等场景中反复出现.
\end{Eg}

\begin{Note}[为什么是二次加速]
    经典概率算法在概率密度(振幅平方)上操作, 一次尝试成功概率为 $p$;
    量子算法在振幅上做旋转, $k$ 次迭代把振幅从 $\sin(\theta/2) \approx \sqrt{p}$
    放大到 $\sin((2k+1)\theta/2)$, 只需 $k = O(1/\sqrt{p})$ 次 —— 这就是
    "量子处理振幅而非概率" 带来的平方级改进.
\end{Note}

\subsection{最优振幅放大 (OAA)}

普通振幅放大的问题是: 若不知道 $p$, 迭代次数 $k$ 难以精确选取, 可能"过冲"或"欠冲",
成功概率达不到 1. 当 $p$ \textbf{已知}时, 可以做到\textbf{确定性}地放大到成功概率恰为 1,
这就是\textbf{最优振幅放大(optimal amplitude amplification, OAA)}.

\begin{Thm}[OAA: 已知概率时的确定性放大]
    在定理 \ref{thm:AA} 的假设下, 若 $p$ 已知, 且可以访问受控反射
    $cR_{\mathrm{good}} = \ket{1}\bra{1} \otimes R_{\mathrm{good}} + \ket{0}\bra{0} \otimes I$,
    则存在量子算法用 $O(1/\sqrt{p})$ 次受控反射查询, 确定性地制备出
    $\ket{0}\ket{\mathrm{good}}$.
\end{Thm}

\begin{Proof}
    令 $\theta = 2\arcsin\sqrt{p}$, 取整数 $k = O(1/\theta) = O(1/\sqrt{p})$,
    并选取 $p' \le p$ 使得 $\theta' := 2\arcsin\sqrt{p'} = \dfrac{\pi}{2k+1}$
    (例如取 $k = \lceil \pi/(2\theta)\rceil$, 则 $2k+1 \ge \pi/\theta$,
    故 $p' = \sin^2(\pi/(2(2k+1))) \le \sin^2(\theta/2) = p$).
    引入一个辅助比特, 施加单比特酉 $B$ 使得
    \[ B\ket{0} = \sqrt{1 - p'/p}\,\ket{0} + \sqrt{p'/p}\,\ket{1}, \]
    于是新初态 $\ket{\phi_0'} = (B \otimes U_0)\ket{0}\ket{0^n}$ 对新好空间
    $\Pi'_{\mathrm{good}} = \ket{1}\bra{1} \otimes \Pi_{\mathrm{good}}$ 的成功振幅恰为
    $\sqrt{p'}$, 即 $\theta' = \pi/(2k+1)$. 此时
    \[ \big|\bra{\mathrm{good}'} G'^{k} \ket{\phi_0'}\big|
       = \sin\!\Big((2k+1)\frac{\theta'}{2}\Big) = \sin\frac{\pi}{2} = 1, \]
    迭代 $k$ 次后\textbf{必然}落在好空间内. 由于 $\Pi'_{\mathrm{good}}$ 由辅助比特
    $\ket{1}$ 标志, 反射为受控反射 $cR_{\mathrm{good}}$, 每步恰用 1 次查询.
    最后对辅助比特施加 $X$ 即得 $\ket{0}\ket{\mathrm{good}}$.
\end{Proof}

\begin{Note}[OAA 的意义与局限]
    OAA 把"成功概率 $\Omega(1)$"提升为"成功概率 $= 1$", 消除了重复试验的开销,
    在 Grover 搜索等 $p$ 已知的场景中尤为有用. 代价是需要受控且可逆的 oracle.
    当 $p$ 未知时, 经典策略是随机选取迭代次数 $m \in [1, 2^r)$ 后检验,
    成功概率在 $2^r > \pi/\theta$ 时平均仍有常数; 现代算法则用 \textbf{固定点振幅放大}
    (fixed-point amplitude amplification, 基于 QSVT)在未知 $p$ 的情况下逼近 100% 成功,
    这将在下一模块中介绍.
\end{Note}

\section{Lie--Trotter 与哈密顿量模拟}

\subsection{问题设定}

\begin{Def}[哈密顿量模拟]
    给定厄米矩阵 $H$(哈密顿量)与初态 $\ket{\psi(0)}$, 目标是制备
    \begin{equation}
        \ket{\psi(t)} = e^{-iHt}\ket{\psi(0)}.
    \end{equation}
    这是 Schrödinger 方程 $i\partial_t \ket{\psi(t)} = H\ket{\psi(t)}$ 的解.
\end{Def}

\begin{Note}
    哈密顿量模拟是 Feynman 1982 年提出量子计算机设想的直接动机,
    也是 QPE 及其应用(基态能量估计、HHL 等)中 $U = e^{-iHt}$ 的实现途径.
    与一般线性 ODE 不同, 由于 $e^{-iHt}$ 是酉算子, 我们无需存储态的完整历史,
    只需关注 $t$ 时刻的状态. 本模块介绍最经典的方法 —— \textbf{乘积公式(product formula)},
    即 Lie--Trotter 分解; 更先进的方法(block-encoding、qubitization、QSVT)见下一模块.
\end{Note}

\subsection{Lie 乘积公式与一阶 Trotter}

当 $H = H_1 + H_2$ 且 $[H_1, H_2] \neq 0$ 时, $e^{-itH} \neq e^{-itH_1}e^{-itH_2}$,
但\textbf{Lie 乘积公式}给出

\begin{equation}
    e^{-itH} = \lim_{L \to \infty} \Big(e^{-i(t/L)H_1} e^{-i(t/L)H_2}\Big)^{L}.
\end{equation}

记 $\tau = t/L$ 为步长, 一步的\textbf{局部误差}为 $O(\tau^2)$:

\begin{Thm}[一阶 Trotter 公式]
    \begin{equation}
        \norm{e^{-i\tau H} - e^{-i\tau H_1} e^{-i\tau H_2}} = O(\tau^2),
    \end{equation}
    其中常数与 $\norm{H_1}$, $\norm{H_2}$ 有关. 对总时间 $t$ 做 $L$ 步, 全局误差
    \begin{equation}
        \norm{e^{-itH} - \big(e^{-i\tau H_1} e^{-i\tau H_2}\big)^{L}} = O(L\tau^2)
        = O\!\Big(\frac{t^2}{L}\Big),
    \end{equation}
    因此达到精度 $\varepsilon$ 需要
    \begin{equation}
        L = O\!\Big(\frac{t^2}{\varepsilon}\Big)
    \end{equation}
    步, 每步调用 $e^{-i\tau H_1}$, $e^{-i\tau H_2}$ 各一次.
\end{Thm}

\subsection{交换子型误差界}

上面的 $O(\tau^2)$ 界中的常数可以精确刻画. 对一阶 Trotter, 记 $U(\tau) = e^{-i\tau H_1}e^{-i\tau H_2}$,
用 Module1 中的 \textbf{Duhamel 原理}分析: $U(\tau)$ 满足
$i\partial_\tau U = HU + e^{-i\tau H_1}[H_2, e^{i\tau H_1}]\cdots$ 形式的含源方程,
解出

\begin{equation}
    \norm{e^{-i\tau H} - e^{-i\tau H_1} e^{-i\tau H_2}}
    \le \frac{\tau^2}{2}\norm{[H_1, H_2]} \le (\tau \alpha)^2,
\end{equation}

其中 $\alpha = \max\{\norm{H_1}, \norm{H_2}\}$, 第二式用到了
\begin{equation}
    \norm{[H_1,H_2]} \le 2\norm{H_1}\,\norm{H_2}.
\end{equation}
于是达到精度 $\varepsilon$ 所需的全局步数为
$L = O\big(t^2\norm{[H_1,H_2]}/\varepsilon\big)$.

\begin{Eg}[横向场 Ising 模型: 交换子界优于算子范数界]
    一维横向场 Ising 模型
    \[ H = -J\sum_{i=1}^{n-1} Z_i Z_{i+1} - g\sum_{i=1}^{n} X_i =: H_1 + H_2. \]
    $H_1$ 内部各项两两对易, $H_2$ 内部亦然, 故 $e^{-i\tau H_1}$, $e^{-i\tau H_2}$ 可精确实现
    (乘积分解为局部旋转). 但 $\norm{H_1} = O(n)$, $\norm{H_2} = O(n)$, 算子范数界给出
    $\alpha^2 = O(n^2)$, 即 $L = O(n^2 t^2/\varepsilon)$; 而
    $[Z_i Z_{i+1}, X_k] \neq 0$ 仅当 $k = i$ 或 $k = i+1$, 故
    $\norm{[H_1, H_2]} = O(n)$, 交换子界给出更紧的 $L = O(n t^2/\varepsilon)$.
    这说明: 误差界应尽量用交换子范数, 而不是简单的算子范数.
\end{Eg}

\begin{Note}[向量范数界]
    算子范数界对最坏情形的初态成立; 若关心特定初态 $\ket{\psi_0}$, 误差可按
    \begin{equation}
        \frac{t^2}{2}\,\max_{0 \le s \le t}\norm{[H_1,H_2]\ket{\psi(s)}}
    \end{equation}
    估计, 往往比算子范数界小得多(例如势场问题中可与系统规模无关).
\end{Note}

\subsection{对称分裂与高阶 Suzuki 公式}

\begin{Thm}[二阶 Trotter(对称分裂, Strang 分裂)]
    \begin{equation}
        \norm{e^{-i\tau H} - e^{-i\tau H_1/2} e^{-i\tau H_2} e^{-i\tau H_1/2}} = O(\tau^3),
    \end{equation}
    达到精度 $\varepsilon$ 需要 $L = O(t^{3/2} \varepsilon^{-1/2})$ 步, 每步 3 次局部演化.
\end{Thm}

\begin{Thm}[高阶 Suzuki 公式]
    高阶公式可以通过\textbf{组合方法(composition)}递归构造. 若 $U_p$ 是 $p$ 阶公式
    (局部误差 $O(\tau^{p+1})$), 且系数满足 $\sum_j s_j = 1$, $\sum_j s_j^{p+1} = 0$,
    则组合 $U_p(s_m\tau)\cdots U_p(s_1\tau)$ 至少是 $p+1$ 阶.
    特别地, \textbf{Trotter--Suzuki 递推}给出 $2k$ 阶公式
    \begin{equation}
        U_{2k}(\tau) = U_{2k-2}(s_k\tau)\; U_{2k-2}\big((1-4s_k)\tau\big)\; U_{2k-2}(s_k\tau),
        \qquad s_k = \frac{1}{4 - 4^{1/(2k-1)}},
    \end{equation}
    每个 $U_{2k}$ 调用 5 次 $U_{2k-2}$, 故 $U_{2k}$ 含至多 $5^{k-1}$ 次 $U_2$.
    一般地, $p$ 阶公式达到精度 $\varepsilon$ 需要
    \begin{equation}
        L = O\!\big(t^{1+1/p}\,\varepsilon^{-1/p}\big)
    \end{equation}
    步; 当 $p \to \infty$ 时渐近为 $L = O(t^{1+o(1)}\varepsilon^{-o(1)})$.
\end{Thm}

\begin{Note}[阶数与代价的权衡]
    高阶公式改进了对 $t$ 和 $\varepsilon$ 的依赖($1/p$ 指数), 但每步的指数个数随阶数
    指数增长($5^{k-1}$), 实践中 2 至 6 阶往往最优. 常系数 $s_k$ 的构造
    (Yoshida--Suzuki 三跳法与 5 级 Trotter--Suzuki 法)是数值分析中组合方法的直接移植.
\end{Note}

\subsection{局部演化的实现与交换情形}

乘积公式的每个因子 $e^{-i\tau H_j}$ 仍需分解为基本门. 两种常见情形:

\begin{Eg}[交换哈密顿量的精确模拟]
    若 $H = \sum_j H_j$ 且各 $H_j$ 两两对易, 则 $e^{-iHt} = \prod_j e^{-iH_j t}$ 精确成立,
    无需任何分裂误差. 例如 TFIM 的 $H_1 = \sum_i Z_i Z_{i+1}$ 与
    $H_2 = \sum_i X_i$ 各自内部对易.
\end{Eg}

\begin{Eg}[两比特 Pauli 项的电路]
    对 $H = ZZ$(省略常系数), 由 $X e^{-itZ} X = e^{itZ}$ 得
    \begin{equation*}
    \Qcircuit @C=1em @R=.8em {
        \lstick{\ket{x}} & \ctrl{1} & \qw & \ctrl{1} & \rstick{\ket{x}} \qw \\
        \lstick{\ket{y}} & \targ & \gate{R_z(2t)} & \targ & \rstick{\ket{y}} \qw
    }
    \end{equation*}
    即 $e^{-it\,Z\otimes Z} = \mathrm{CNOT}\cdot\big(I \otimes R_z(2t)\big)\cdot\mathrm{CNOT}$.
    一般 $k$ 比特 Pauli 项可先用 CNOT 链计算奇偶性, 施加单个 $R_z$ 后再用 CNOT 链还原
    (共 $2(k-1)$ 个 CNOT 与 1 个 $R_z$); 含 $X$, $Y$ 的项则先用 Clifford 共轭
    ($X = HZH$, $Y = SHS^{\dagger} Z S H S^{\dagger}$ 类关系)对角化后再处理.
\end{Eg}

\end{document}